\documentclass[english,11pt,a4paper]{article}
\usepackage[margin=1in]{geometry}
\usepackage{fontspec}
\setmainfont{lmroman10-regular.otf}[BoldFont=lmroman10-bold.otf,ItalicFont=lmroman10-italic.otf,BoldItalicFont=lmroman10-bolditalic.otf]
\setsansfont{lmsans10-regular.otf}[BoldFont=lmsans10-bold.otf]
\setmonofont{lmmono10-regular.otf}[Scale=MatchLowercase]
\newfontfamily\LeanSymbols{seguisym.ttf}[Path=typeset_runtime/]
\newlength{\LeanCharWidth}
\usepackage{amsmath,amssymb}
\usepackage[english]{babel}
\usepackage{microtype,parskip}
\usepackage{graphicx,float,caption}
\usepackage{longtable,booktabs,array,calc}
\usepackage{xcolor}
\usepackage{xurl}
\usepackage{hyperref,bookmark}
\usepackage{etoolbox,needspace}
\setcounter{secnumdepth}{-\maxdimen}
\setlength{\emergencystretch}{3em}
\providecommand{\tightlist}{\setlength{\itemsep}{0pt}\setlength{\parskip}{0pt}}
\newcounter{none}
\AtBeginEnvironment{longtable}{\Needspace{8\baselineskip}\small}
\pretocmd{\section}{\Needspace{5\baselineskip}}{}{}
\pretocmd{\subsection}{\Needspace{4\baselineskip}}{}{}
\pretocmd{\subsubsection}{\Needspace{4\baselineskip}}{}{}
\captionsetup{font=small,labelfont=bf}
\urlstyle{same}
\providecommand{\pandocbounded}[1]{#1}
\hypersetup{pdftitle={Paper I: formalization and reproducibility},pdfauthor={Juan Pablo Traverso Gianini},pdflang={en-US},colorlinks=true,linkcolor=blue,urlcolor=blue,citecolor=blue}
\title{Paper I: formalization and reproducibility}
\author{Juan Pablo Traverso Gianini\\\small Independent researcher, Santiago, Chile\\\small \href{mailto:jtraverso@gmail.com}{jtraverso@gmail.com}\\\small ORCID: 0009-0003-6068-4096}
\date{}
\begin{document}
\raggedbottom
\maketitle
\textbf{Paper I · Formal companion v1.4}\\
19 October 2026

This document accompanies the mathematical article. It preserves the
formal declarations, scope qualifications and computational evidence
previously included in the article. The source freeze and its recorded
verification are unchanged. No Lean/Lake build was run for this edition.
Historical section and table numbers in the records below refer to the
preceding article; they are retained as locators, not renumbered as new
verification results. Original references {[}1{]}, {[}2{]}, etc. are
listed at the end.

\href{https://github.com/jtraverso/erdos-81-chordal-clique-partitions/tree/eb335845bfa23375d5040ea31fc4f4b366346287/preprints/PAPER_I/05_formalization}{Frozen
sources and historical evidence / Fuentes congeladas y evidencias
históricas} ·
\href{https://github.com/jtraverso/erdos-81-chordal-clique-partitions/tree/eb335845bfa23375d5040ea31fc4f4b366346287/preprints/PAPER_I/01_manuscript}{Previous
article / Artículo anterior}

\subsubsection{Theorem 1.1}\label{theorem-1.1}

\textbf{Formalization status.} The Paper I Lean v1.2 freeze uses Lean 4
/ Mathlib v4.28.0 {[}6,7{]} and records Theorem 1.1 as
\nolinkurl{paperI_main_sharp}. The declaration \nolinkurl{paperI_main}
supplies a weaker companion form with additive term \(n\). In these
recorded names, \nolinkurl{sharp} refers to the quadratic coefficient
\(1/6\), not to optimality of the additive term. The axiom gate covers
the main theorem, \nolinkurl{PaperI.assembly_sharp},
\nolinkurl{PaperI.Split.residual_duality}, the finite packing--covering
duality of Proposition B.1, and the interface declarations listed in
Section 10. Their recorded footprint is \nolinkurl{propext},
\nolinkurl{Classical.choice}, and \nolinkurl{Quot.sound}, with no
project-level mathematical axiom. Section 10 gives the freeze and
archive record.

\subsection{10. Reproducibility and proof
record}\label{reproducibility-and-proof-record}

The proof is analytic. Computational experiments were used during
development to search for counterexamples and detect transcription
errors, but no computation is a logical premise of the theorem. For that
reason, computational regression material is kept outside the paper
rather than included as a second appendix.

The accompanying calculation ledger records:

\begin{itemize}
\tightlist
\item
  every exact identity;
\item
  every inequality direction;
\item
  the complete loss accounting;
\item
  the branch \(q=0\);
\item
  the orbit regimes;
\item
  the final conversion to \(n\).
\end{itemize}

The author reports an internal AI-assisted adversarial audit of the
theorem statement, quantitative bound, and proof assembly. This is a
validation record, not a claim of external peer review.

\subsubsection{Formalization status in
Lean}\label{formalization-status-in-lean}

The snapshot at \nolinkurl{PAPER_I/05_formalization/lean_v1.2_freeze}
contains the Lean sources and a successful recorded 8,034-job build
against Lean 4 / Mathlib v4.28.0 {[}6,7{]}. It contains
\nolinkurl{paperI_main_sharp}, whose statement is
\text{\(\Phi\le n^2/6+n/2\)}, together with the weaker companion
declaration \nolinkurl{paperI_main} and the interface results listed
below. Its axiom gate also queries \nolinkurl{PaperI.assembly_sharp} and
\nolinkurl{PaperI.Split.residual_duality}. Appendix C records the main
declaration. The exact archive was subsequently reproduced in an
external clean room with the same 8,034-job result and the same
theorem-level axiom footprint.

The frozen axiom report records only Lean's standard foundational
axioms, \nolinkurl{propext}, \nolinkurl{Classical.choice}, and
\nolinkurl{Quot.sound}, for the main theorem and the listed declarations;
it records no project-level mathematical axiom. These are logical
foundations used throughout Mathlib, not additional mathematical
assumptions of the paper.

The release artifact also records the finite-dimensional
packing-covering duality used in Proposition B.1. In Lean it appears as
\nolinkurl{residual_duality}, an application of the general theorem
\nolinkurl{covering_packing_duality}. The artifact records that theorem
as derived from a finite Farkas lemma and Mathlib's closed-cone
infrastructure, rather than postulated as an external axiom. Schrijver
{[}4, Chapter 7{]} remains the classical reference for the written
proof.

The frozen source archive is \nolinkurl{PAPER_I_lean_v1.2_freeze.zip},
with SHA-256
\nolinkurl{0181506408644fc1f8872d711de5a98a500f4052aa295bcd6f8c82776694fd3a}.
It is stored under
\nolinkurl{PAPER_I/05_formalization/lean_v1.2_freeze} and excludes
project caches and compiled objects. The freeze build target includes
several \nolinkurl{Contrib} modules used for possible upstreaming; that
contribution lane is a byproduct and is not part of the paper's trusted
mathematical assumptions. The axiom report above controls the logical
dependency statement. The external clean-room reproduction applies to
these exact archive bytes; the manuscript revision does not alter the
formal artifact.

\begin{center}\rule{0.5\linewidth}{0.5pt}\end{center}

\subsection{Appendix C. Formalized
statement}\label{appendix-c.-formalized-statement}

The accompanying formal artifact records the following Lean 4
declaration for Theorem 1.1:

\begin{quote}
{\LeanSymbols\fontsize{8}{10}\selectfont
\noindent \nolinkurl{theorem}\hspace{0.45em}\nolinkurl{paperI_main_sharp}\hspace{0.45em}\nolinkurl{(G}\hspace{0.45em}\nolinkurl{:}\hspace{0.45em}\nolinkurl{Split)}\hspace{0.45em}\nolinkurl{:}\par
\noindent \hspace{1.8em}\nolinkurl{G.Phi}\hspace{0.45em}\nolinkurl{≤}\hspace{0.45em}\nolinkurl{(G.n}\hspace{0.45em}\nolinkurl{:}\hspace{0.45em}\nolinkurl{ℝ)}\hspace{0.45em}\nolinkurl{^}\hspace{0.45em}\nolinkurl{2}\hspace{0.45em}\nolinkurl{/}\hspace{0.45em}\nolinkurl{6}\hspace{0.45em}\nolinkurl{+}\hspace{0.45em}\nolinkurl{(G.n}\hspace{0.45em}\nolinkurl{:}\hspace{0.45em}\nolinkurl{ℝ)}\hspace{0.45em}\nolinkurl{/}\hspace{0.45em}\nolinkurl{2}\par
\noindent \par
\noindent \nolinkurl{--}\hspace{0.45em}\nolinkurl{weaker}\hspace{0.45em}\nolinkurl{companion}\hspace{0.45em}\nolinkurl{form}\hspace{0.45em}\nolinkurl{(larger}\hspace{0.45em}\nolinkurl{additive}\hspace{0.45em}\nolinkurl{term):}\par
\noindent \nolinkurl{theorem}\hspace{0.45em}\nolinkurl{paperI_main}\hspace{0.45em}\nolinkurl{(G}\hspace{0.45em}\nolinkurl{:}\hspace{0.45em}\nolinkurl{Split)}\hspace{0.45em}\nolinkurl{:}\par
\noindent \hspace{1.8em}\nolinkurl{G.Phi}\hspace{0.45em}\nolinkurl{≤}\hspace{0.45em}\nolinkurl{(G.n}\hspace{0.45em}\nolinkurl{:}\hspace{0.45em}\nolinkurl{ℝ)}\hspace{0.45em}\nolinkurl{^}\hspace{0.45em}\nolinkurl{2}\hspace{0.45em}\nolinkurl{/}\hspace{0.45em}\nolinkurl{6}\hspace{0.45em}\nolinkurl{+}\hspace{0.45em}\nolinkurl{(G.n}\hspace{0.45em}\nolinkurl{:}\hspace{0.45em}\nolinkurl{ℝ)}\par
}
\end{quote}

Here \nolinkurl{Split} packages the data of a split graph used in the
proof: a clique of size \nolinkurl{p}, an independent vertex type, and the
neighborhoods of the independent vertices inside the clique. The field
\nolinkurl{n} is the total number of vertices. The quantity \nolinkurl{Phi} is
defined as \texttt{\textbar{}E\textbar{}\ −\ 2·nu3star}, where
\nolinkurl{nu3star} is the supremum of the feasible fractional
triangle-packing values. Elsewhere in the series this fractional deficit
is written \(\Phi^*\), while the unstarred \(\Phi\) is reserved for the
integral deficit \(|E(G)|-2\nu_3(G)\); the Lean identifier is retained
to match the recorded declaration. With these definitions,
\nolinkurl{paperI_main_sharp} is precisely

\[
|E(G)|-2\nu_3^*(G)\le \frac{n^2}{6}+\frac{n}{2},
\]

the inequality stated in Theorem 1.1 (declaration
\nolinkurl{paperI_main_sharp}). The companion declaration
\nolinkurl{paperI_main} records the weaker additive term \(n\).

The supplied build report records zero \nolinkurl{sorry}. The supplied
axiom report is:

\begin{quote}
{\LeanSymbols\fontsize{8}{10}\selectfont
\noindent \nolinkurl{PaperI.Split.paperI_main_sharp}\hspace{0.45em}\nolinkurl{depends}\hspace{0.45em}\nolinkurl{on}\hspace{0.45em}\nolinkurl{axioms:}\hspace{0.45em}\nolinkurl{[propext,}\hspace{0.45em}\nolinkurl{Classical.choice,}\hspace{0.45em}\nolinkurl{Quot.sound]}\par
\noindent \nolinkurl{PaperI.paperI_main}\hspace{4.05em}\nolinkurl{depends}\hspace{0.45em}\nolinkurl{on}\hspace{0.45em}\nolinkurl{axioms:}\hspace{0.45em}\nolinkurl{[propext,}\hspace{0.45em}\nolinkurl{Classical.choice,}\hspace{0.45em}\nolinkurl{Quot.sound]}\par
\noindent \nolinkurl{PaperI.assembly_sharp}\hspace{2.7em}\nolinkurl{depends}\hspace{0.45em}\nolinkurl{on}\hspace{0.45em}\nolinkurl{axioms:}\hspace{0.45em}\nolinkurl{[propext,}\hspace{0.45em}\nolinkurl{Classical.choice,}\hspace{0.45em}\nolinkurl{Quot.sound]}\par
\noindent \nolinkurl{PaperI.Split.residual_duality}\hspace{0.45em}\nolinkurl{depends}\hspace{0.45em}\nolinkurl{on}\hspace{0.45em}\nolinkurl{axioms:}\hspace{0.45em}\nolinkurl{[propext,}\hspace{0.45em}\nolinkurl{Classical.choice,}\hspace{0.45em}\nolinkurl{Quot.sound]}\par
}
\end{quote}

These are the three standard foundational axioms of Lean and Mathlib. No
project-specific axiom is used.

\textbf{Table 2. Frozen formalization perimeter (Lean v1.2 freeze).} The
frozen Lean 4 / Mathlib v4.28.0 snapshot records the proved core with
only Lean's standard foundational axioms and no project-level
mathematical axiom. An external clean-room build and theorem-level axiom
check reproduced the exact archive.

{\def\LTcaptype{none} % do not increment counter
\begin{longtable}[]{@{}
  >{\raggedright\arraybackslash}p{(\linewidth - 4\tabcolsep) * \real{0.3333}}
  >{\raggedright\arraybackslash}p{(\linewidth - 4\tabcolsep) * \real{0.3333}}
  >{\raggedright\arraybackslash}p{(\linewidth - 4\tabcolsep) * \real{0.3333}}@{}}
\toprule\noalign{}
\begin{minipage}[b]{\linewidth}\raggedright
Node
\end{minipage} & \begin{minipage}[b]{\linewidth}\raggedright
Statement
\end{minipage} & \begin{minipage}[b]{\linewidth}\raggedright
Axiom footprint
\end{minipage} \\
\midrule\noalign{}
\endhead
\bottomrule\noalign{}
\endlastfoot
\nolinkurl{paperI_main_sharp} (Theorem 1.1) &
\(\lvert E(G)\rvert-2\nu_3^*(G)\le n^2/6+n/2\) & \nolinkurl{propext},
\nolinkurl{Classical.choice}, \nolinkurl{Quot.sound} \\
\nolinkurl{paperI_main} (weaker companion bound) &
\(\lvert E(G)\rvert-2\nu_3^*(G)\le n^2/6+n\) & \nolinkurl{propext},
\nolinkurl{Classical.choice}, \nolinkurl{Quot.sound} \\
\nolinkurl{lpVal_concave}, \nolinkurl{lpVal_ge_min} (§5 interface) &
abstract concavity / pure-profile lower bound & \nolinkurl{propext},
\nolinkurl{Classical.choice}, \nolinkurl{Quot.sound} \\
\nolinkurl{lpVal_lipschitz} (§5 interface) & parametric stability of the
fractional value & \nolinkurl{propext}, \nolinkurl{Classical.choice},
\nolinkurl{Quot.sound} \\
\nolinkurl{residual_duality} (Prop. B.1) & finite packing--covering
duality & \nolinkurl{propext}, \nolinkurl{Classical.choice},
\nolinkurl{Quot.sound} \\
\end{longtable}
}

\subsubsection{Proposition B.1}\label{proposition-b.1}

\textbf{Formalization note.} The v1.2 freeze records Proposition B.1 as
the Lean declaration \nolinkurl{PaperI.Split.residual_duality}, an
instance of the general finite strong-duality theorem
\nolinkurl{covering_packing_duality}. The artifact records the latter as
proved from a finite Farkas lemma, using the closedness of a finitely
generated convex cone and the hyperplane-separation theorem for closed
convex cones. Thus the duality used here is derived inside the formal
development, not postulated as a project-specific axiom. Schrijver
remains the classical reference for the written argument. The
corresponding sources, build log, and axiom report are delivered in
\nolinkurl{PAPER_I/05_formalization/lean_v1.2_freeze}.

\subsubsection{1. Introduction}\label{introduction}

\textbf{Series notation.} In cross-paper discussion, the fractional
deficit of Paper I is written \(\Phi^*(G)=|E(G)|-2\nu_3^*(G)\). Paper II
writes the same quantity as \(\Phi_\tau(G)=|E(G)|-2\tau_3^*(G)\), while
Paper III reserves the unstarred \(\Phi(G)\) for the integral deficit
\(|E(G)|-2\nu_3(G)\). The Lean identifier \nolinkurl{Phi} in the Paper I
artifact is retained for traceability.

\subsubsection{Remark (fixed-polytope
concavity)}\label{remark-fixed-polytope-concavity}

The reduction just used is an instance of a purely convex-analytic fact,
stated here in abstract form. Fix a feasible polytope that does not
depend on the profile and an objective that depends affinely on the
profile mass \(\mu=(\mu_S)\). Because the optimal value of a linear
program over a fixed polytope is concave in its objective vector, the
value at a mixed profile dominates the corresponding convex combination
of pure-profile values, and hence their minimum: \[
M_b(\kappa_\mu)\ \ge\ \sum_S \tfrac{\mu_S}{q}\,M_b(\kappa^{S,q})\ \ge\ \min_S M_b(\kappa^{S,q}).
\] This is recorded in the formal development as \nolinkurl{lpVal_concave}
and \nolinkurl{lpVal_ge_min}. A companion estimate,
\nolinkurl{lpVal_lipschitz}, bounds the change in the fractional value
when the baseline or the profile distribution is perturbed. It is
intended for robust and near-split variants. Both are stated for a fixed
feasible polytope with an affinely varying objective; neither requires
symmetry, which enters only afterwards, on a single pure profile, in
Section 6.

\subsection{References}\label{references}

{[}1{]} P. Erdős, E. T. Ordman, and Y. Zalcstein, ``Clique partitions of
chordal graphs,'' \emph{Combinatorics, Probability and Computing}
\textbf{2} (1993), no. 4, 409--415.

{[}2{]} G.-T. Chen, P. Erdős, and E. T. Ordman, ``Clique partitions of
split graphs,'' in \emph{Combinatorics, Graph Theory, Algorithms and
Applications} (Beijing, 1993), World Scientific, 1994, pp.~21--30,
author-hosted scan:
\url{https://ordman.net/MathResearch/CEOClique_Parts.pdf}, accessed
August 21, 2026.

{[}3{]} T. F. Bloom, ``Erdős Problem \#81,'' \emph{Erdős Problems},
\url{https://www.erdosproblems.com/81}, accessed August 21, 2026.

{[}4{]} A. Schrijver, \emph{Theory of Linear and Integer Programming},
Wiley-Interscience Series in Discrete Mathematics and Optimization, John
Wiley \& Sons, Chichester, 1986.

{[}5{]} M. S. Cavers, \emph{Clique Partitions and Coverings of Graphs},
M.Math. essay, University of Waterloo, 2005.

{[}6{]} L. de Moura and S. Ullrich, ``The Lean 4 Theorem Prover and
Programming Language,'' in \emph{Automated Deduction -- CADE 28},
Lecture Notes in Computer Science \textbf{12699}, Springer, 2021,
pp.~625--635, \url{https://doi.org/10.1007/978-3-030-79876-5_37}.

{[}7{]} The mathlib Community, ``The Lean Mathematical Library,'' in
\emph{Proceedings of the 9th ACM SIGPLAN International Conference on
Certified Programs and Proofs (CPP 2020)}, ACM, 2020, pp.~367--381,
\url{https://doi.org/10.1145/3372885.3373824}.

\end{document}
